\section{Conclusion}
\label{sec:conclusion}

Capacity utilization has been treated as something an economy reports.
This chapter has treated it as something an economy produces. Once the
denominator of the ratio is recognized as the outcome of a
transformation---accumulated capital becoming usable productive capacity
through a choice of technique made under distributive
conflict---utilization stops being a primitive and becomes a residual of that
transformation. The elasticity $\theta$ is what has to be recovered
first, and everything the ratio can be made to say depends on recovering
it without assuming its value.

The apparatus for that recovery came from two frameworks that had not
previously been joined at this point. Okishio's account of decentralized
accumulation establishes that private investment commitments generate
disproportionality as a structural norm rather than as a friction around
proportional reproduction \citep{Okishio1961, Okishio2022}. Vidal's
account of the internally divided firm establishes that the organization
converting capital into operating capacity cannot maximize control and
cooperation at once, and settles instead into arrangements that sustain
production without resolving the conflict inside it
\citep{Vidal2022, Vidal2019}. Together they rule out any automatic
passage from accumulation to capacity, which is what makes $\theta$ an
object with a determinate value worth estimating rather than a parameter
to be normalized to one.

The estimates say two things about the center. American capacity
formation has stayed strictly below the Harrodian knife-edge across
1931--2024 under both the homogeneous and the heterogeneous capital
specification: accumulation persistently outruns the capacity it creates.
And the elasticity has drifted upward across the postwar period, from
$0.371$ in 1960 to $0.429$ in 2024, for a reason that does not vindicate
the accumulation regime that produced it. Effective wage pressure fell
from $0.698$ to $0.504$ over those decades, partly because the wage share
declined and partly because investment migrated into intellectual
property, which now exceeds a quarter of fixed capital. A rising
transformation elasticity coincided with falling utilization, which stood
at $93.8\%$ in 2024 after a trough of $74.1\%$ in 1982. Capital became
somewhat better at building capacity precisely as the distributive
settlement that would have realized that capacity was dismantled.

The periphery shows the same relation cut short from outside. For Chile
over 1943--2010, a composite balance-of-payments index built from reserve
cover and the import-financing gap identifies a threshold at
$\hat{\gamma} = 0.425$ by profile likelihood ratio search
\citep{Hansen2000}. In the unconstrained state a rising wage share
induces mechanization as cost-minimization predicts, with an estimated
interaction of $+3.29$. When the threshold binds, a triple interaction of
$-3.15$ removes nearly all of that response. Chilean capitalists under
foreign-exchange rationing keep the incentive to substitute machinery for
labor and lose the means, because the machinery is imported. The Unidad
Popular period is the clearest case: a twenty-point rise in the gross
wage share to $\omega_{1972} = 0.6903$ produced almost no mechanization
response, and utilization fell to $0.72$ by 1975.

What follows for how utilization is used. Any empirical strategy that
recovers utilization by fitting output to capital has imposed
$\theta = 1$ before measuring anything, and the estimates here put the
value near $0.4$ for the United States and between $0.68$ and $0.95$ for
Chile. Utilization series built that way do not measure the distance
between output and capacity; they measure the residual of an assumption.
For the long-running debate over whether firms converge on a normal
utilization rate, the finding is more direct still: the denominator that
debate presupposes is itself a function of the distributive conflict the
debate treats as separate from it.

Two extensions follow with enough specificity to be worth naming. The
Chilean series terminate in 2010, which places the unwinding of the
commodity super-cycle outside the sample; splicing the historical
reserve, terms-of-trade, and money series to contemporary Central Bank
sources through 2024 would test whether the external gate has migrated
again, from markup expansion toward lithium-era rent capture, and the
threshold specification estimated here transfers to that extended sample
without modification. Separately, the two economies are identified on
different specifications, which limits what the comparison can carry;
constructing an intellectual-property capital series for Chile, or an
import-biased mechanization gap for the United States, would let a single
specification be estimated on both panels and would turn the
sign-structure comparison offered here into a quantitative one.

The scope of this chapter is deliberately bounded. It does not present a
complete crisis theory, a fully specified model of profitability
dynamics, or an exhaustive dependency framework. It reconstructs capacity
utilization as a structurally mediated object and identifies the
transformation elasticity as the hinge of its empirical recovery in the
center and the periphery, and it leaves the broader crisis and
profitability architecture to the chapters that follow.
